\section{Introduction}

\label{sec:intro}

The architecture of modern applications is undergoing a fundamental shift. AI-native applications---those that incorporate LLM inference, embedding computation, retrieval-augmented generation, and model serving as core components---have distinct infrastructure requirements that existing PaaS offerings do not address.

\subsection{The AI Infrastructure Gap}

Traditional PaaS platforms were designed for request-response web applications with predictable resource consumption. AI applications differ in several critical ways:

\begin{enumerate}[leftmargin=1.4em]
    \item \textbf{Heterogeneous compute}: AI workloads require GPUs for inference, CPUs for pre/post-processing, and specialized hardware (TPUs, Inferentia) for cost optimization. Traditional PaaS provides only CPU-based containers.
    \item \textbf{Model lifecycle}: Models must be downloaded (often multi-GB), loaded into GPU memory, warmed up, and version-managed. Traditional PaaS treats all application artifacts as stateless code bundles.
    \item \textbf{Token-based economics}: LLM inference is priced per-token rather than per-request, requiring fundamentally different cost modeling and optimization strategies.
    \item \textbf{Streaming responses}: LLM applications typically stream responses via Server-Sent Events or WebSockets, requiring persistent connections that conflict with traditional HTTP load balancers.
    \item \textbf{Vector storage}: RAG applications require vector databases co-located with inference for low-latency retrieval, a primitive not provided by traditional PaaS.
\end{enumerate}

\subsection{Contributions}

\begin{enumerate}[leftmargin=1.4em]
    \item An AI-aware build system with automatic model dependency detection and GPU-optimized container images (\S\ref{sec:build}).
    \item A unified resource model treating AI primitives as first-class platform resources (\S\ref{sec:resources}).
    \item A cost-optimized autoscaler for inference workloads (\S\ref{sec:autoscaler}).
    \item Comprehensive evaluation against competing platforms (\S\ref{sec:evaluation}).
    \item Production deployment analysis from 14 months of operation (\S\ref{sec:production}).
\end{enumerate}

